\begin{definition}[7.8.2]
\EGAReaderTarget{ega4.c04.def.7-8-2}{IV.2.s7.8.2-fr}
Se dice que un anillo $A$ es \textup{excelente} si es noetheriano y satisface las
condiciones siguientes:
\begin{enumerate}
\item[\emph{(i)}] $A$ es universalmente catenario \textup{(o, lo que equivale a lo mismo
(5.6.3, (i)),} para todo ideal primo $\mathfrak p$ de $A$,
$A_{\mathfrak p}$ es universalmente catenario).
\item[\emph{(ii)}] Para todo ideal primo $\mathfrak p$ de $A$, las fibras
formales de $A_{\mathfrak p}$ son geométricamente regulares.
\item[\emph{(iii)}] Para todo anillo cociente $B$ de $A$ que sea un dominio de
integridad y toda extensión finita puramente inseparable $K'$ del cuerpo de
fracciones $K$ de $B$, existe una sub-$B$-álgebra finita $B'$ de $K'$,
que contiene $B$, tiene $K'$ como cuerpo de fracciones y es tal que el conjunto
de puntos regulares de $\operatorname{Spec}(B')$ contiene un abierto no vacío.
\end{enumerate}

\upshape
Con esta terminología, lo esencial de los resultados del \S~7 y de una parte
del \S~6 se resume como sigue:
\end{definition}
